
% Place here your custom commands, it will be included in presentation.tex

\newcommand{\normalurl}[1]{\href{#1}{#1}}

\newcommand{\footnoteurl}[1]{\footnote{\scriptsize\url{#1}}}

\newcommand{\rocketbox}[1]{%
	\awesomebox[circlRed]{1pt}{\faRocket}{circlBlue}{#1}%
}


% =========================================================
% STORY STATUS FOOTER
% =========================================================
%
% Usage:
%
% 1. Put this in your preamble
%
% 2. Before each section, define the current phase:
%
%    \storyphase{Incident}
%    \storyphase{Correlation}
%    \storyphase{Monitoring}
%    \storyphase{Operations}
%    \storyphase{Governance}
%
% 3. The footer automatically highlights the current phase
%
% =========================================================

% ---------------------------------------------------------
% Current phase holder
% ---------------------------------------------------------

\definecolor{phaseactive}{RGB}{255,255,255}
\definecolor{phaseinactive}{RGB}{220,220,220}

\newcommand{\currentstoryphase}{}
\newcommand{\allstoryphases}{}

\newcommand{\storyphase}[1]{%
    \gdef\currentstoryphase{#1}%
    \hypertarget{#1}{}
}

\newcommand{\setphases}[1]{%
    \gdef\allstoryphases{#1}%
}

\newcommand{\phasetag}[1]{%
    \expandafter\ifstrequal\expandafter{\currentstoryphase}{#1}
        {\color{phaseactive}\textbf{\scriptsize \hyperlink{#1}{#1}}}
        {\color{phaseinactive}\hyperlink{#1}{#1}}%
}

% ---------------------------------------------------------
% Render one phase
% ---------------------------------------------------------
% \newcommand{\phasetag}[1]{%
%     \expandafter\ifstrequal\expandafter{\currentstoryphase}{#1}
%         {\color{phaseactive}\textbf{#1}}
%         {\color{phaseinactive}#1}%
%     }


% ---------------------------------------------------------
% Render all phases
% ---------------------------------------------------------
% \newcommand{\renderphases}{%
%     \foreach \x [count=\i] in \allstoryphases {%
%         \phasetag{\x}%
%         \ifnum\i<\foreachcount
%             \hspace{1.2em}|\hspace{1.2em}%
%         \fi }%
%     }

\newif\iffirstphase

\newcommand{\renderphases}{%
    \global\firstphasetrue
    \foreach \x in \allstoryphases {%
        \iffirstphase
            \global\firstphasefalse
        \else
            \hspace{1.2em}|\hspace{1.2em}%
        \fi
        \expandafter\phasetag\expandafter{\x}%
    }%
}


% ---------------------------------------------------------
% Custom footer
% ---------------------------------------------------------

\setbeamertemplate{footline}
{
    \leavevmode%
    \begin{beamercolorbox}[
        wd=\paperwidth,
        ht=3.2ex,
        dp=1.2ex,
        leftskip=1em,
        rightskip=1em
    ]{author in head/foot}
    \tiny
    \insertframenumber{} / \inserttotalframenumber
    \hfill
    \renderphases
    \end{beamercolorbox}
}

% \setbeamertemplate{footline}
% {
%     \leavevmode%
%     \begin{beamercolorbox}[
%         wd=\paperwidth,
%         ht=3.2ex,
%         dp=1.2ex,
%         leftskip=1em,
%         rightskip=1em
%     ]{author in head/foot}
%     \tiny
%     \insertframenumber{} / \inserttotalframenumber
%     \hfill
%     \phasetag{Incident}
%     \hspace{1.2em}|\hspace{1.2em}
%     \phasetag{Correlation}
%     \hspace{1.2em}|\hspace{1.2em}
%     \phasetag{Monitoring}
%     \hspace{1.2em}|\hspace{1.2em}
%     \phasetag{Operations}
%     \hspace{1.2em}|\hspace{1.2em}
%     \phasetag{Governance}
%     \hspace{1.2em}
%     \end{beamercolorbox}
% }

% Optional: remove navigation symbols
\setbeamertemplate{navigation symbols}{}

% Nice badge.
% Usage:
% \entitytag{border-color}{fill-color}{text-label}
% \entitytag{blue!60!black}{blue!5}{Internal MISP}
\newcommand{\entitytag}[3]{%
\tikz[baseline=(X.base)]{
    \node[
        draw=#1,
        fill=#2,
        rounded corners=2pt,
        inner xsep=1.5pt,
        inner ysep=0pt
    ] (X) {\strut #3};
}}


% Create a link button in the bottom-right corner that jump to the target
% Usage:
% \jumpto{hyperlinkID}{textlabel}
\newcommand{\jumpto}[2]{%
\begin{tikzpicture}[remember picture,overlay]
    \node[anchor=south east]
        at ([yshift=10pt]current page.south east)
    {
        \hyperlink{#1}{\beamergotobutton{#2}}
    };
    \end{tikzpicture}
}
% Create a link button in the bottom-right corner that return to the target
% Usage:
% \returnto{hyperlinkID}{textlabel}
\newcommand{\returnto}[2]{%
\begin{tikzpicture}[remember picture,overlay]
    \node[anchor=south east]
        at ([yshift=0pt]current page.south east)
    {
        \hyperlink{#1}{\beamerreturnbutton{#2}}
    };
    \end{tikzpicture}
}